\documentclass[11pt,a4paper]{article}

\usepackage[utf8]{inputenc}
\usepackage[T1]{fontenc}
\usepackage[margin=2.5cm]{geometry}
\usepackage{booktabs}
\usepackage{amsmath}
\usepackage{xcolor}
\usepackage{hyperref}
\hypersetup{colorlinks=true, linkcolor=blue!50!black, urlcolor=blue!50!black,
            citecolor=blue!50!black}

\title{Classifying hate and counter speech in Bluesky replies to news outlets:\\
       a proposal grounded in measurement}
\author{BlueX project}
\date{14 August 2026 --- revised 21 August 2026}

\begin{document}
\maketitle

\begin{abstract}
\noindent
This proposal sets out how the BlueX corpus---2.42\,million replies to five news
outlets on Bluesky---can be labelled and classified for hate speech, and later
counter speech. Every design choice is justified by a measurement already taken on
this corpus, and several otherwise plausible approaches are excluded because they
were tried and failed. The central finding is that off-the-shelf toxicity models do
not detect hate: once incivility is controlled for they are \emph{anti-correlated}
with it. The central constraint is that the distant-supervision strategy which
succeeded in earlier work on German Twitter is unavailable here, and the
platform-moderation labels that might have replaced it turn out to be almost
entirely retracted. We propose a low-cost architecture---paragraph embeddings
trained unsupervised on the full corpus, a small supervised classifier on top, an
abstention threshold, and score-banded human validation---and state in advance what
would falsify it. The classifier is organised as a \emph{consensus committee} of
deliberately decorrelated local models; the committee's disagreement, together with
a sample of its unanimous verdicts, directs where the scarce human labelling budget
is spent.
\end{abstract}

\section{Motivation and scope}

BlueX studies hate speech, counter speech and disinformation in replies to news
outlets, and how these change over time and around public events. This document
covers the \emph{classification} problem for hate. Counter speech is deferred until
hate can be identified, because counter speech is relational---a reply is counter
speech relative to what it answers---and the practical route to it runs through
identifying hate first. Disinformation is discussed only to record that it currently
has no viable data path (\S\ref{sec:disinfo}).

\section{The corpus}

Collection is complete in the sense that matters: every tracked outlet has been
scraped back to its first post on the platform, and every reply tree has been walked.

\begin{table}[h]
\centering
\begin{tabular}{lrrl}
\toprule
Outlet & Root posts & Zero-reply roots & First post \\
\midrule
spiegel.de              & 86{,}444 & 48\,\% & 2023-10-02 \\
theguardian.com         & 57{,}303 & 12\,\% & 2024-11-15 \\
zeit.de                 & 43{,}647 & 64\,\% & 2023-10-23 \\
nytimes.com             & 40{,}499 &  3\,\% & 2023-06-22 \\
tagesschau.bsky.social  & 40{,}306 & 25\,\% & 2023-09-14 \\
\bottomrule
\end{tabular}
\caption{Corpus composition. Total: 2{,}423{,}185 posts, 268{,}199 roots,
267{,}496 distinct reply authors. Reply trees: 253{,}166 walked to completion; zero
replies whose root post is missing.}
\end{table}

Roughly 82\,\% of replies are English and 14\,\% German. The wide variation in
zero-reply rate---3\,\% for the New York Times, 64\,\% for ZEIT---means that
``replies per outlet'' and ``replies per engaged post'' are different quantities.
Any analysis must state which it uses.

Two structural limitations are unrecoverable. First, replies deleted before
collection are gone: comparing a June~2026 archive against a re-scrape of the same
27{,}961 roots showed a $-1.08\,\%$ loss over two months, with 17.1\,\% of threads
losing at least one reply. Because moderation removes the most extreme content
first, this erosion is biased \emph{against} the material of interest. Second, when
an account is removed its content becomes unreachable from both the AppView and the
account's own personal data server; only the identity record survives.

\section{Label inventory, corrected}
\label{sec:labels}

Bluesky's moderation labels were harvested in full for posts (1{,}382{,}554
subjects swept) and accounts (205{,}391 swept), with zero failed batches. The
harvest records the \texttt{neg} flag, which marks a label as \emph{negated}---that
is, retracted. Accounting for negation changes the picture substantially, and an
earlier internal summary of this project's own data failed to do so.

\begin{table}[h]
\centering
\begin{tabular}{lrr}
\toprule
Account label & Active & Negated \\
\midrule
\texttt{!suspend}     &   43 & 4{,}903 \\
\texttt{!takedown}    &   20 & 1{,}757 \\
\texttt{needs-review} & 2{,}865 & 5{,}127 \\
\bottomrule
\end{tabular}
\hspace{1cm}
\begin{tabular}{lrr}
\toprule
Post label & Active & Negated \\
\midrule
\texttt{rude}        & 881 &  2 \\
\texttt{spam}        & 371 & 12 \\
\texttt{intolerant}  & 178 &  1 \\
\texttt{threat}      &  60 &  1 \\
\bottomrule
\end{tabular}
\caption{Active versus negated moderation labels. Only 63 accounts currently carry
an active suspension or takedown---not the 6{,}723 obtained by counting
applications and retractions together.}
\end{table}

The usable supervision is therefore \textbf{238 active hate-relevant post labels}
(\texttt{intolerant}, \texttt{threat}, \texttt{extremist},
\texttt{intolerant-race}) and \textbf{881 \texttt{rude}} labels. The post labels are
sound: only 37 of 1{,}771 are negated, and negated labels were excluded from every
experiment reported below.

\paragraph{Consequence.} Account-level distant supervision is not available. The 63
actioned accounts account for 463 replies in total. This is the single most
important difference between BlueX and the prior work discussed next.

\section{What the prior work did, and why it does not transfer}

Garland, Ghazi-Zahedi, Young, H\'ebert-Dufresne and Galesic classified hate and
counter speech on German Twitter at scale~\cite{garland2020countering}, and analysed
the resulting dynamics over four years~\cite{garland2020impact}. Four components of
that method transfer very differently.

\paragraph{Distant supervision from self-labelling groups (does not transfer).}
Germany in 2017--2018 hosted two organised movements that identified themselves:
\emph{Reconquista Germanica}, a coordinated far-right group, and \emph{Reconquista
Internet}, a counter-speech movement founded in opposition. Millions of tweets from
their members served as naturally labelled training data. BlueX's corpus---replies
to news outlets---contains no known equivalent, and as \S\ref{sec:labels} shows, the
platform's own moderation labels cannot substitute: 463 replies is not a training
set.

\paragraph{Ensemble of paragraph embeddings with regularised logistic regression
(transfers well).} The published classifier paired document embeddings with
regularised logistic regression in an ensemble. This is a low-cost, CPU-trainable
design and forms the basis of \S\ref{sec:proposal}.

\paragraph{Abstention via a confidence threshold, validated against neutral
timelines (transfers, with a caveat).} The diagnostics plot recall and
percentage-classified against a threshold~$\tau$, separately for known hate
accounts, known counter accounts, and \emph{neutral} timelines. As $\tau$ rises the
fraction of neutral content classified falls towards zero: the model declines to
judge material it is unsure about. This is a principled answer to the ``neither''
problem.

The caveat is arithmetic. Abstention controls the false-positive \emph{rate}, not
the false-positive \emph{count}, and the count depends on the base rate. Reading
$\tau = 0.9$ from the published curves, roughly 5\,\% of neutral content is still
classified while about 20\,\% of hate-account content is classified at 0.88 recall.
In a corpus that is 95\,\% neutral and 1\,\% hate, that yields on the order of
twenty-five false positives per true positive. The prior corpus consisted of
political conversations seeded by two organised movements, where the base rate of
hate was far higher. \textbf{BlueX's base rate has never been measured, and it
determines whether abstention alone suffices.}

\paragraph{Crowdsourced validation stratified by model score (transfers directly).}
Validation samples were drawn in bands across the classifier's score range
($0.5$--$0.55$, $0.55$--$0.6$, \ldots, $0.95$--$1.0$) for both classes. This yields
calibration---what a score of $0.7$ actually means---rather than a single aggregate
figure, and it supports principled threshold selection. We adopt it directly
(\S\ref{sec:validation}).

\paragraph{A note on reported performance.} The published macro $F_1$ of
$0.76$--$0.97$ is on \emph{balanced} out-of-sample test sets. This is the same
qualifier our own experiments encountered (\S\ref{sec:measurements}), and is worth
stating plainly rather than discovering later.

\section{What has been measured on this corpus}
\label{sec:measurements}

All experiments use one evaluation set: 235 active hate-labelled replies
(\emph{positive}), 872 \texttt{rude} replies (\emph{hard negative}), and 940
randomly drawn unlabelled replies (\emph{easy negative}).

\subsection{Sentiment carries no signal}

Apple's \texttt{NLTagger} sentiment score separates hate from random replies at
AUC~$0.508$---chance. The medians of the two distributions are identical
($-0.600$). Replies to news accounts are already overwhelmingly negative (78\,\% of
controls score below zero), so negativity is the corpus baseline rather than a
discriminating feature.

\subsection{Toxicity models detect incivility, and are anti-correlated with hate}

\begin{table}[h]
\centering
\begin{tabular}{lrrr}
\toprule
Detector head & hate vs.\ random & hate vs.\ \texttt{rude} & \texttt{rude} vs.\ random \\
\midrule
\texttt{unbiased-toxic-roberta\#identity\_attack} & 0.901 & 0.518 & 0.919 \\
\texttt{unbiased-toxic-roberta\#toxicity}         & 0.846 & \textbf{0.198} & \textbf{0.946} \\
\texttt{textdetox\#toxic}                         & 0.829 & \textbf{0.188} & 0.944 \\
\texttt{twitter-roberta-base-hate\#hate}          & 0.814 & 0.533 & --- \\
keyword lexicon (baseline)                        & 0.511 & 0.508 & --- \\
\bottomrule
\end{tabular}
\caption{ROC AUC by task. An AUC of $0.198$ means the model rates a
moderator-labelled \emph{rude} post as more toxic than a moderator-labelled
\emph{hate} post 80\,\% of the time.}
\end{table}

Toxicity models score profanity, insult and obscenity. The label
\texttt{intolerant}---Bluesky's term for discrimination against protected
groups---is frequently expressed in clean language. The models measure a different
construct, accurately. \textbf{Any pipeline that filters by incivility before
detecting hate will systematically discard the cleanly-worded hate that matters
most.}

A validated \emph{incivility} detector falls out as a by-product: the
\texttt{toxicity} head separates \texttt{rude} from random at $0.946$, and the whole
corpus has been scored with it (2{,}085{,}088 replies, zero failures).

\subsection{Hate is learnable, but only against the wrong background}

Fine-tuning directly on the hate-versus-\texttt{rude} distinction, with stratified
5-fold cross-validation:

\begin{table}[h]
\centering
\begin{tabular}{lrr}
\toprule
Model & hate vs.\ \texttt{rude} (CV) & hate vs.\ random \\
\midrule
\texttt{cardiffnlp/twitter-roberta-base} (fine-tuned) & $0.959 \pm 0.014$ & 0.61--0.68 \\
\texttt{xlm-roberta-base} (fine-tuned)                & $0.932 \pm 0.033$ & 0.61--0.68 \\
TF--IDF $+$ logistic regression                        & $0.943 \pm 0.019$ & --- \\
keyword lexicon                                        & $0.508 \pm 0.011$ & --- \\
\bottomrule
\end{tabular}
\caption{Cross-validated performance. Train/test gaps were 0.03--0.05, so these are
not memorisation.}
\end{table}

Two conclusions follow. First, the distinction \emph{is} learnable: the off-the-shelf
failures reflected a mismatch of construct, not an absence of signal. Second, and
decisively for deployment, a model trained on hate versus \texttt{rude} answers
``given one of these two, which is it?'' rather than ``is this hate?''. Against the
true corpus, which is overwhelmingly neither, performance falls to $0.61$--$0.68$.

Note also that TF--IDF with logistic regression lies within one standard deviation of
a fine-tuned transformer. \textbf{Any proposal to use a large model must beat
TF--IDF in the same harness, or justify itself on generalisation rather than
headline accuracy.}

\section{Proposal}
\label{sec:proposal}

\subsection{Architecture}

\paragraph{Stage 0: measure the base rate.} Draw a uniform random sample of replies
and label it by hand. This yields the prevalence of hate in the deployment
distribution---currently unknown, and the number that determines whether any
threshold-based approach is viable (\S4). It comes first because it can falsify
everything built on top of it.

\paragraph{Stage 1: unsupervised representation.} Train paragraph embeddings
(doc2vec, or a comparable sentence encoder) \emph{unsupervised} on all 2.15\,M
replies. No labels are required. This is the key economy: the representation is
learned from the deployment distribution itself, ordinary content included, which is
precisely the blind spot that produced the fall to $0.61$--$0.68$. It also learns
the corpus's own bilingual vocabulary, which off-the-shelf models lack. Cost:
CPU-hours, no GPU.

\paragraph{Stage 2: small supervised classifier, organised as a committee.} Fit
regularised logistic regression on the embeddings using the 238 active hate labels,
the 881 \texttt{rude} labels, and negatives drawn from the Stage~0 sample. Following
\cite{garland2020countering}, use an ensemble rather than a single model and retain
at least one interpretable member: TF--IDF's strong showing suggests lexical
features carry much of the signal, and coefficients can be read directly. The
ensemble is structured as a consensus committee whose members are chosen for
decorrelation, not merely for accuracy (\S\ref{sec:committee}).

\paragraph{Stage 3: abstention.} Classify only where confidence exceeds $\tau$;
otherwise return ``unjudged''. Choose $\tau$ from the calibration curves of
\S\ref{sec:validation} rather than by convention.

\paragraph{Stage 4: full-corpus inference.} Encoder-class inference runs at
$\sim$145 posts per second on project hardware, so the full corpus takes about
4.5 hours---demonstrated end to end by the completed incivility pass.

\subsection{Why doc2vec rather than a larger model}

Three reasons, each measured rather than assumed. TF--IDF already reaches $0.943$,
so shallow lexical representations capture most of the available signal. The
supervised label budget is roughly one thousand examples, far below what fine-tuning
a large model can exploit without overfitting. And an embedding trained on the corpus
itself represents the majority class---ordinary replies---whereas a model fine-tuned
only on hate versus \texttt{rude} has never seen it.

\subsection{The consensus committee, and disagreement as a sampling signal}
\label{sec:committee}

\paragraph{Composition.} A committee's votes carry information only insofar as its
members err \emph{independently}; three fine-tunes of the same encoder would agree
confidently for identical reasons. The members are therefore chosen to represent
different families of failure:

\begin{enumerate}
\item \textbf{TF--IDF $+$ logistic regression} --- pure surface-lexical signal
      ($0.943$ in the benchmark), near-zero inference cost, coefficients readable.
\item \textbf{doc2vec $+$ logistic regression} --- distributional but order-blind,
      trained unsupervised on the corpus itself (Stage~1), so it alone has seen the
      majority class.
\item \textbf{A fine-tuned contextual encoder} --- \texttt{twitter-roberta-base} or
      \texttt{xlm-roberta-base}, the only member sensitive to word order and
      context; $\sim$145 posts per second on project hardware.
\item \textbf{Optionally, a small local LLM (8B class) as tie-breaker.} Roughly two
      orders of magnitude slower than the encoder, it never scores the full corpus;
      it is consulted only on the disagreement region the cheap members define. Its
      verdicts are automatic labels and fall under the audit obligation of
      \S\ref{sec:validation}, item~3.
\end{enumerate}

Apple's \texttt{NLTagger} is excluded on measurement: AUC $0.508$
(\S\ref{sec:measurements}).

\paragraph{Aggregation.} Members return calibrated scores, not hard votes. The
committee's verdict is the mean score; its \emph{disagreement} is the spread
(standard deviation) across members. Spread is a finer signal than vote counts: a
2--1 split at scores $(0.51, 0.49, 0.52)$ is near-consensus, while $(0.95, 0.05,
0.50)$ at the same majority is maximal conflict. Abstention (Stage~3) applies to the
mean; high spread is itself grounds to abstain, since a committee that cannot agree
should not be counted.

\paragraph{Disagreement-driven human review.} The committee's output partitions the
corpus into three regions, each of which earns human labels for a different reason:

\begin{itemize}
\item \textbf{High spread (the split region).} Query-by-committee sampling: labels
      here are maximally informative for improving the members, and at a base rate
      of a few percent they are worth far more than uniform draws, most of which
      land on unambiguous neutral content.
\item \textbf{Unanimous positive.} All current members inherit the same construct
      bias---they measure incivility more readily than hate (the $0.198$
      dissociation)---so the committee can be \emph{confidently wrong in a
      correlated direction}: a rude-but-not-hateful post drawing three assured
      ``hate'' scores. Reviewing a sample of unanimous positives measures precision
      exactly where correlated error hides, and coincides with the top band of the
      score-banded calibration already adopted.
\item \textbf{Unanimous negative.} A smaller sample bounds the false-negative rate
      at the confident end; without it, ``the committee found nothing'' is
      unfalsifiable.
\end{itemize}

\paragraph{Statistical discipline.} Labels drawn by disagreement are maximally
non-random. They must never enter the base-rate estimate, and if they are used to
retrain the members that selected them, they cease to be held-out gold. Both traps
are handled by machinery that already exists: every label records its sampling
frame, so committee-drawn labels carry a distinct frame kind (recording the member
models, their versions, and the score snapshot that triggered selection) and are
partitionable from uniform-random labels in every later analysis. Stage~0's uniform
sample remains the only source of prevalence estimates. Committee-drawn labels are
spent on \emph{validation and calibration} of the committee; any retraining on them
converts them explicitly from gold to training data, recorded as such, and they are
thereafter excluded from evaluation.

\paragraph{Ordering.} The committee is built \emph{after} Stage~0, not before. Only
the measured base rate says what the disagreement region contains: at 1\,\%
prevalence versus 5\,\%, the same split region has a very different composition,
and the number of disagreement cases worth reviewing changes accordingly.

\section{Labelling programme}
\label{sec:validation}

Human annotation capacity is a few hundred items, one annotator. It is therefore
spent only where nothing cheaper can substitute.

\begin{enumerate}
\item \textbf{Base-rate estimation} (Stage~0). A uniform random sample, hand-labelled.
      Nothing else provides an unbiased prevalence estimate.
\item \textbf{Score-banded calibration}, following \cite{garland2020countering}.
      Sample uniformly across the classifier's score range and label each band,
      yielding precision per band, a calibration curve, and a defensible $\tau$ from
      a few hundred labels.
\item \textbf{Audit of any automatic labels.} If an LLM or external API generates
      training labels at volume, a hand-labelled subset must measure how good those
      labels are. Without it the pipeline rests on an unmeasured assumption.
\item \textbf{Committee-directed review} (\S\ref{sec:committee}). Samples from the
      high-disagreement region, the unanimous-positive band, and the
      unanimous-negative band, each recorded under its own sampling frame. These
      labels calibrate and validate the committee; they never feed the base-rate
      estimate.
\end{enumerate}

Human labels are \textbf{held-out gold and never training data}: with a few hundred
available, spending them on training would leave nothing to validate against.

Labels may be drawn from filtered pools where efficiency demands it, but
\textbf{every label records its sampling frame}---the filter, the pool size, the
batch, and the random seed. Without this, labels drawn from a high-incivility filter
become indistinguishable afterwards from a uniform draw, and any prevalence estimate
across them is silently biased.

\paragraph{Reliability.} With a single annotator, Cohen's $\kappa$ is unavailable.
The obtainable substitute is intra-rater agreement: re-label a subset blind after an
interval. This is weaker than inter-rater agreement and must be reported as such,
not quietly omitted.

\paragraph{Presentation.} Annotation shows the reply together with its root post and
immediate parent, since hate reads differently in context; this is always possible,
as every reply's parent is present in the store. Model outputs, incivility scores and
moderation labels are \emph{not fetched} into the annotation view---not merely
hidden---because these labels are the only independent ground truth the project has.

\section{Counter speech}

Deferred until hate can be identified, for two reasons. It is relational: ``that is
disgusting, delete this'' is counter speech under a hateful parent and something else
under a news headline. And it has \emph{no} external labels---the complete sweep of
1{,}382{,}554 replies produced none, because moderators label violations rather than
virtues.

The intended approach, following the structure of \cite{garland2020impact}, is to
identify accounts that habitually engage in counter speech and examine their replies
to hate. This is more tractable than classifying every reply for a relational
property, and it yields a natural denominator: of threads containing hate, how many
drew a response from a known counter-speaker?

\section{Disinformation}
\label{sec:disinfo}

The complete post-label sweep produced three relevant labels in 1.38\,M replies
(\texttt{misleading}: 2, \texttt{rumor}: 1). There is no external signal and no data
path. A scope decision is required: either commit to annotation from scratch, or
remove disinformation from the research questions. Retaining it as an unstated
aspiration is the worst of the three options.

\section{Risks, stated in advance}

\begin{itemize}
\item \textbf{The base rate may be too low for abstention to suffice.} This is the
      principal risk; Stage~0 exists to detect it before anything is built on top.
\item \textbf{Automatic labels may be poor and unaudited.} Any use of LLM or API
      labelling requires a hand-labelled audit subset.
\item \textbf{Evaluating on the convenient task.} Hate versus \texttt{rude} gives
      $0.96$; hate versus random gives $0.61$--$0.68$. The second is the deployment
      task and must be the headline.
\item \textbf{Reintroducing an incivility filter.} Toxicity is anti-correlated with
      hate here. The result is counter-intuitive and therefore easy to forget.
\item \textbf{Correlated committee confidence.} Every member inherits the same
      incivility-versus-hate construct bias, so unanimity is not independence.
      Consensus alone must never be promoted to ground truth; the unanimous bands
      are sampled for human review precisely because agreement can be confidently
      wrong in the same direction.
\item \textbf{Treating moderation labels as ground truth for recall.} They record
      what was reported and actioned. A model scoring poorly against them may be
      finding hate nobody reported.
\item \textbf{Single-annotator reliability.} Intra-rater agreement is a weaker
      substitute for $\kappa$ and must be described honestly.
\end{itemize}

\begin{thebibliography}{9}
\bibitem{garland2020countering}
J.~Garland, K.~Ghazi-Zahedi, J.-G.~Young, L.~H\'ebert-Dufresne, and M.~Galesic.
\newblock Countering hate on social media: large scale classification of hate and
counter speech.
\newblock \emph{arXiv:2006.01974}, 2020.
\newblock \url{https://arxiv.org/abs/2006.01974}

\bibitem{garland2020impact}
J.~Garland, K.~Ghazi-Zahedi, J.-G.~Young, L.~H\'ebert-Dufresne, and M.~Galesic.
\newblock Impact and dynamics of hate and counter speech online.
\newblock \emph{arXiv:2009.08392}, 2020.
\newblock \url{https://arxiv.org/abs/2009.08392}
\end{thebibliography}

\end{document}
